% Generated draft: requires mathematical and source review.
\documentclass[11pt]{article}
\usepackage[margin=25mm]{geometry}
\usepackage{fontspec}
\setmainfont{Latin Modern Roman}
\usepackage{amsmath,amssymb,mathtools,unicode-math}
\setmathfont{Latin Modern Math}
\usepackage{enumitem,needspace,xurl,hyperref}
\setlength{\parindent}{0pt}
\setlength{\parskip}{5pt}
\setlength{\emergencystretch}{3em}
\newcommand{\N}{\mathbb{N}}
\newcommand{\Z}{\mathbb{Z}}
\newcommand{\Q}{\mathbb{Q}}
\newcommand{\R}{\mathbb{R}}
\newcommand{\C}{\mathbb{C}}
\newcommand{\F}{\mathbb{F}}
\newcommand{\A}{\mathbb{A}}
\newcommand{\PP}{\mathbb{P}}
\newcommand{\E}{\mathbb{E}}
\newcommand{\eps}{\varepsilon}
\newcommand{\dd}{\,\mathrm{d}}
\begin{document}
\section*{UnsolvedMath ID: 10400053}
\textbf{Source problem code:} catalogue entry
\subsection{Definitions and mathematical statement}
Let \(K\subset S^3\) be an oriented knot, considered up to ambient isotopy, and let \(O\) be the oriented unknot. Let \(Z(K)\) be the Kontsevich invariant in a fixed convention, valued in the degree completion of the appropriate chord/Jacobi-diagram algebra. Its degree completion retains every finite diagram degree; write \(Z_d(K)\) for degree \(d\).

Does there exist a nontrivial oriented knot \(K\), meaning one not isotopic to \(O\), such that
\[
 Z(K)=Z(O),\qquad\text{equivalently }\quad
 Z_d(K)=Z_d(O)\ \text{for every }d\ge0?
\]
This is precisely a request for a counterexample to unknot detection by the full Kontsevich invariant. It is not the weaker problem of matching the unknot through one fixed finite degree, nor the stronger assertion that \(Z\) distinguishes every pair of knots.

The source does not reproduce its diagram/coefficient and framing convention. One must use the same normalization and, if a framed version is used, a matching specified framing for both knots. In particular \(Z(O)\) is retained as written and is not silently replaced by the algebra unit, since that depends on normalization. The source points to Problem 2.6 as equivalent, but that problem is not supplied; no extra equivalence is inferred from its number.

The summary identifies this as the open unknot-detection question, and supplies no modern resolution or exact auxiliary citation. The 2002 problem and later open label remain source assessments pending current primary-literature review.
\subsection{Short English statement}
Can a nontrivial oriented knot have exactly the same full Kontsevich invariant as the unknot? Equality is required in every degree, using matching normalizations and framing conventions.
\subsection{Sources}
\begin{itemize}
\item Local registry: \url{data/unsolvedmath/problems.json}, permanent UnsolvedMath ID \(10400053\), source code catalogue entry; supplied statement, background, and research summary, updated 2026-08-06.
\item Ohtsuki, \emph{Problems on invariants of knots and 3-manifolds} (2002), Problem 3.3, referring to Problem 2.6, PDF page 53: \url{https://msp.org/gtm/2002/04/gtm-2002-04-024s.pdf}. Supplied PDF-font-structure extraction accessed 2026-07-29; provenance distinguishes 231 numbered items.
\item The record supplies no separate current resolution reference or full text of Problem 2.6.
\end{itemize}
\end{document}
